\section{CRT candidate and terminal residue}

This chapter records the CRT layer of the unified core.  The layer
reduces the final valuation inequality to exact modular arithmetic and
two size conditions.  The mathematical derivation is complete; the
Lean statement \texttt{unified\_core\_final\_no\_hge} that assembles
the last step remains pending.

\subsection{Setup}

Let $j$ be the reset depth, let $s$ be the terminal depth, and write
\begin{equation*}
  n=s-j,\qquad \Delta=W_s-W_j.
\end{equation*}
Let $L$ be the length of the block tail and let $B$ be the block
molecule of the tail word.  The tail is a valid word from the block
head $r_{j,0}$ to the terminal state $r_s$, so the exact word-orbit
identity gives
\begin{equation*}
  2^{\Delta}\,r_s=5^n\,r_{j,0}+B.
\end{equation*}

\subsection{The terminal odd part}

The terminal odd part is
\begin{equation*}
  w_{\mathrm{term}}=\frac{3r_s+1}{2^{L+4}}.
\end{equation*}
Multiplying the tail equation by $3\cdot2^{\Delta}$ and adding
$2^{\Delta}$ gives
\begin{equation*}
  2^{\Delta+L+4}\,w_{\mathrm{term}}
    =2^{\Delta}(3r_s+1)
    =3\cdot5^n\,r_{j,0}+3B+2^{\Delta}.
\end{equation*}
This identity is the exact algebra connecting the orbit equation to
the valuation inequality.  It is the identity recorded by
\texttt{rj0\_\allowbreak spec\_\allowbreak eq}.

\subsection{The valuation obstruction}

The target inequality is
\begin{equation*}
  \vtwo{5^{L+3}w_{\mathrm{term}}+1}\le H_s-2.
\end{equation*}
Its negation is
\begin{equation*}
  \vtwo{5^{L+3}w_{\mathrm{term}}+1}\ge H_s-1,
\end{equation*}
which holds exactly when
\begin{equation*}
  5^{L+3}w_{\mathrm{term}}+1\equiv0\pmod{2^{H_s-1}}.
\end{equation*}
Since $5$ is odd, $5^{-(L+3)}$ exists modulo $2^{H_s-1}$, and the
negation is exactly
\begin{equation*}
  w_{\mathrm{term}}\equiv-5^{-(L+3)}
  \pmod{2^{H_s-1}}.
\end{equation*}
Let $u$ be the least nonnegative residue of $-5^{-(L+3)}$ modulo
$2^{H_s-1}$.

\subsection{The CRT equation}

Assume, toward the contradiction, that the valuation inequality fails.
Then $w_{\mathrm{term}}\equiv u\pmod{2^{H_s-1}}$, so
\begin{equation*}
  w_{\mathrm{term}}=u+m'\,2^{H_s-1}
\end{equation*}
for some integer $m'\ge0$.  Substituting into the terminal identity
gives the exact CRT equation
\begin{equation*}
  3\cdot5^n\,r_{j,0}+3B+2^{\Delta}
    =2^{\Delta+L+4}\left(u+m'\,2^{H_s-1}\right).
\end{equation*}
The CRT candidate $r_{j,0}$ is the least nonnegative solution of this
equation.  The equation is linear in $r_{j,0}$, so the candidate is
well defined; its explicit form is used in
\texttt{rj0\_\allowbreak spec\_\allowbreak eq}.

\subsection{The terminal residue $y^*$}

The residue
\begin{equation*}
  y^*=
  -\left(w_{\mathrm{term}}+5^{-(L+3)}\right)
  \left(3\cdot5^s\right)^{-1}
  \pmod{2^{H_s-1}}
\end{equation*}
is the normalised failure residue.  Because $3\cdot5^s$ is odd, its
inverse modulo $2^{H_s-1}$ exists.  By construction,
\begin{equation*}
  y^*=0\iff w_{\mathrm{term}}\equiv-5^{-(L+3)}
  \pmod{2^{H_s-1}},
\end{equation*}
so the valuation inequality is equivalent to $y^*\ne0$.  The Lean
statement is
\texttt{terminal\_\allowbreak bound\_\allowbreak iff\_\allowbreak
yStar\_\allowbreak ne\_\allowbreak zero}, compiled.

\subsection{The two size conditions}

The first size condition controls the correction term:
\begin{equation*}
  3B+2^{\Delta}\le3\cdot5^n-2\cdot4^n.
\end{equation*}
The second controls the leading term:
\begin{equation*}
  3\cdot5^s+\left(3\cdot5^n-2\cdot4^n\right)
    \le2^{\Delta+L+4}\,u.
\end{equation*}
Both are consequences of the no-$H_{\ge}$ premises; the Lean statement
\texttt{blockB\_\allowbreak bound\_\allowbreak of\_\allowbreak
no\_\allowbreak hge} packages the first one.

\subsection{From the size conditions to $r_{j,0}\ge5^j$}

From the CRT equation,
\begin{equation*}
  3\cdot5^n\,r_{j,0}
    =2^{\Delta+L+4}\left(u+m'2^{H_s-1}\right)-3B-2^{\Delta}.
\end{equation*}
The right-hand side is at least
\begin{equation*}
  2^{\Delta+L+4}\,u-\left(3\cdot5^n-2\cdot4^n\right),
\end{equation*}
using the first size condition.  The second size condition bounds this
below by $3\cdot5^s$.  Therefore
\begin{equation*}
  3\cdot5^n\,r_{j,0}\ge3\cdot5^s,
\end{equation*}
so
\begin{equation*}
  r_{j,0}\ge5^{s-n}=5^j.
\end{equation*}
This is the content of
\texttt{rj0\_\allowbreak ge\_\allowbreak of\_\allowbreak
size\_\allowbreak conditions\_\allowbreak no\_\allowbreak hge}, compiled.

\subsection{Role of the block-head interval}

The no-$H_{\ge}$ premises also carry interval constraints on the block
head.  The final unified-core statement must compare
$r_{j,0}\ge5^j$ with those constraints.  The comparison is not an
independent analytic estimate: it is the remaining assembly step
\texttt{unified\_\allowbreak core\_\allowbreak final\_\allowbreak
no\_\allowbreak hge}, which is the single pending \texttt{sorry} in the
unified-core audit.

\subsection{Worked modular checks}

The $d=2$ exclusion uses two modular facts.  Modulo $17$,
\begin{equation*}
  5^3=125\equiv6\pmod{17},
\end{equation*}
and the order of $5$ modulo $17$ is $16$, so the congruence
$5^{j-1}\equiv6$ fixes $j-1\equiv3\pmod{16}$.  Modulo $256$,
\begin{equation*}
  5^7=78125\equiv45\pmod{256},
\end{equation*}
and the order of $5$ modulo $256$ is $64$, so $5^{j-1}\equiv45$
fixes $j-1\equiv7\pmod{64}$.  The two congruences contradict each
other modulo $16$.

The $d=3$ family has denominator $109$.  The consistency check is
carried out by repeated squaring modulo $109$:
\begin{center}
\small
\begin{tabular}{@{}rrr@{}}
\toprule
$k$ & $5^k\bmod 109$ & note \\
\midrule
1 & 5 & \\
2 & 25 & \\
4 & 80 & $625-5\cdot109$ \\
8 & 78 & $6400-58\cdot109$ \\
16 & 89 & $6084-55\cdot109$ \\
32 & 73 & $7921-72\cdot109$ \\
64 & 97 & $5329-48\cdot109$ \\
128 & 35 & $9409-86\cdot109$ \\
256 & 26 & $1225-11\cdot109$ \\
512 & 22 & $676-6\cdot109$ \\
\bottomrule
\end{tabular}
\end{center}
Since
\begin{equation*}
  923=512+256+128+16+8+2+1,
\end{equation*}
the product is
\begin{equation*}
  22\cdot26\cdot35\cdot89\cdot78\cdot25\cdot5
  \equiv27\cdot35\cdot89\cdot78\cdot25\cdot5
  \equiv73\pmod{109}.
\end{equation*}
Hence $5^{923}\equiv73\pmod{109}$, and
\begin{equation*}
  8\cdot73=584\equiv39\pmod{109},
\end{equation*}
which is exactly the divisibility required by
\begin{equation*}
  g=\frac{8\cdot5^{923}-39}{109}
\end{equation*}
for the representative $j-1=923$ of the residue class
$j-1\equiv923\pmod{1728}$.

\subsection{Lean statements}

\begin{itemize}[leftmargin=*]
\item \texttt{All36\_\allowbreak 20PremisesNoHge}: the record of the
  no-$H_{\ge}$ premises.
\item \texttt{blockB\_\allowbreak bound\_\allowbreak of\_\allowbreak
  no\_\allowbreak hge}: the block-tail bound.
\item \texttt{rj0\_\allowbreak spec\_\allowbreak eq}: the CRT equation
  and candidate.
\item \texttt{terminal\_\allowbreak bound\_\allowbreak iff\_\allowbreak
  yStar\_\allowbreak ne\_\allowbreak zero}: the valuation equivalence.
\item \texttt{rj0\_\allowbreak ge\_\allowbreak of\_\allowbreak
  size\_\allowbreak conditions\_\allowbreak no\_\allowbreak hge}: the
  size-condition lower bound.
\item \texttt{unified\_\allowbreak core\_\allowbreak final\_\allowbreak
  no\_\allowbreak hge}: pending.
\end{itemize}

\subsection{Status ledger}

\begin{center}
\small
\begin{tabular}{@{}ll@{}}
\toprule
Statement & Status \\
\midrule
CRT equation and candidate definition & Lean: compiled \\
valuation iff $y^*\ne0$ & Lean: compiled \\
size-condition lower bound $r_{j,0}\ge5^j$ & Lean: compiled \\
assembly to the final valuation inequality & Lean: pending \\
\bottomrule
\end{tabular}
\end{center}

The paper does not claim that the pending assembly is closed.
